\section{Expanded Logical Skeleton}

The proof can be compressed into one implication chain, but the referee burden is to see why no object is lost while the chain is compressed.
The chain is:
\[
\text{finite terminal claim}
\Longrightarrow
\text{same-solution terminal record}
\Longrightarrow
\text{ordered face test}
\Longrightarrow
\text{pass or first-face exit}.
\]
The pass branch contradicts terminality by continuation.
The first-face exit branch denies that the alleged obstruction is an in-class terminal member of the same solution.
The only possible way to evade the proof would be to produce a terminal object of the same solution that is neither a member branch nor a first-face exit branch.
The definitions of Pack, Part, and Field are chosen exactly to make such an object visible.

\begin{lemma}[No hidden fourth object after entry]
Let \(Q\) be a terminal record admitted as a record of the original solution.
If \(Q\) has a positive same-fluid body, the original Navier--Stokes equation participates on that body, and there is a positive-scale continuation readout on it, then \(Q\) is the pass branch.
If any one of those three statements fails in ordered form, \(Q\) is a first-face exit branch.
There is no remaining membership predicate in \(\Owork\).
\end{lemma}

\begin{proof}
The working class is defined by exactly the three structures required for same-solution continuation custody.
The first structure is the material packet.
The second is same-equation participation.
The third is positive-scale field readout.
Once those are fixed, the record either keeps all three or loses the first one it fails to keep.
The proof does not introduce a fourth predicate later.
\end{proof}

The strength of the argument is also its vulnerability.
If a referee can exhibit a genuine terminal obstruction of the original solution that does not need positive material body, same-equation participation, or a positive field readout, then the proof would need another face.
The paper's claim is that a finite Clay terminal obstruction cannot have that form.
It must be an obstruction of the same solution, and the three faces are exactly the same-solution structures that make the obstruction a member-side terminal branch.

\section{Maximality And Restart}

The restart step is standard, but it is the only positive theorem used to contradict terminality.
It is therefore important to keep it clean.
For \(s>5/2\), \(H^s\) is an algebra in dimension three and embeds into \(C^1\) after choosing \(s\) high enough for the classical readout used here.
The local Navier--Stokes map in \(H^s\) supplies existence, uniqueness, and continuation for the same branch.
The proof invokes this only after Pack and Part have identified the object being restarted.

\begin{proposition}[Same-branch restart]
Suppose \(Q\) is an admitted terminal record for the original solution and suppose
\[
  \Pack_Q\wedge\Part_{N,Q}\wedge\exists r>0\,\Field_{N,r,Q}.
\]
Then the restart datum is not a new datum for a neighboring problem.
It is the same solution read on a participating positive packet, and the local theorem extends that solution.
\end{proposition}

\begin{proof}
Pack supplies the material identity of the packet.
Part supplies the original equation and pressure law on the packet.
Field supplies the continuation-level \(H^s\) readout.
The local theorem restarts the Navier--Stokes equation from a state of the same solution.
Uniqueness glues the restarted branch to the preterminal branch.
The extension is therefore an extension of the original solution, not a comparison branch.
\end{proof}

This proposition is the reason the proof can be contrapositive.
It does not need to prove a global estimate everywhere in advance.
It needs to show that a terminal record that remains a member has exactly the readout required for local restart.

\section{Why Entry Is Forced By The Clay Negation}

A finite Clay counterexample cannot be a bare sentence saying that smoothness fails.
It must be a failure of the same solution from the original datum.
The object of failure may be hard to extract, and it may be weak or oscillatory, but it is still asserted to belong to that branch.
The proof uses that assertion.

\begin{lemma}[Counterexample claims are object claims]
Let \(u\) be the maximal smooth branch from \(u_0\) on \([0,T_*)\).
If a counterexample claims that \(u\) cannot continue past \(T_*\), then the obstruction must be expressible through a terminal behavior of \(u(t)\) along \(t\uparrow T_*\).
\end{lemma}

\begin{proof}
The counterexample has no other subject.
It is not a statement about an arbitrary sequence of functions, a limiting model, or a separate profile.
It is a statement about the branch \(u\).
Thus its obstruction must be read from that branch.
\end{proof}

This lemma does not claim compactness in a topology strong enough to solve the problem.
It only prevents the counterexample from avoiding identity altogether.
The CM test can accept a weakly described record.
It then asks whether the record has a positive material packet, same-equation participation, and a positive field readout.
Failure of any of these is a face failure, not a failure to start the test.

\section{Pressure Participation}

The nonlocal pressure is one of the places where a proof can accidentally change the problem.
In Leray form,
\[
  \partial_t u+\mathbb P\nabla\cdot(u\otimes u)=\nu\Delta u.
\]
Localizing this equation creates commutators and pressure tails.
Part requires those terms to remain terms of the original equation.
They may be estimated, split, or assigned to exterior accounting, but they may not disappear.

\begin{proposition}[Pressure cannot be dropped before Part]
Assume Pack survives for a packet \(P_Q\).
If the pressure contribution generated by the original whole-space Leray projection is replaced by a local pressure law or ignored as a boundary artifact, then Part has not been proved.
\end{proposition}

\begin{proof}
The packet may still have material identity, but Part asks whether the original Navier--Stokes equation participates on it.
The pressure law is part of that equation.
Dropping or replacing it changes the equation being tested.
Until the pressure term is returned as an accounted term, \(\Part_{N,Q}\) is absent.
\end{proof}

This is why the whole-space Duhamel split is not decorative.
It is the mechanism that keeps pressure and exterior source terms attached to the same branch while the proof separates harmless parts from the surviving exterior packet.

\section{Commutator Accounting}

Cutoffs are useful because the whole-space proof needs to separate compact and exterior regions.
They also create commutators.
The proof treats commutators as accounting terms of the same equation.
A commutator that vanishes is harmless.
A commutator that survives is not a new equation; it becomes part of the exterior packet being tested.

\begin{lemma}[Cutoff terms do not create a fourth branch]
Let a cutoff split produce annular commutator terms.
If the annular terms vanish in the relevant readout, they do not obstruct continuation.
If they do not vanish, they are included in the exterior-source terminal record and are tested by Pack, Part, and Field.
\end{lemma}

\begin{proof}
The annular terms are artifacts of the chosen partition of the original equation.
They have no independent solution identity.
Their only possible terminal role is as part of the same exterior source record.
That record is then decided by the ordered CM test.
\end{proof}

\section{Exterior \(L^2\) Tightness And Bounded Frequencies}

The exterior \(H^s\) problem is difficult because high derivatives can survive even when energy is tight.
The proof uses energy tightness only to remove bounded frequencies.
This is a precise use and does not claim exterior \(H^s\) tightness directly.

For fixed \(N_0\),
\[
  \|P_{\le N_0}f\|_{H^s}
  \le C_s N_0^s\|f\|_{L^2}.
\]
Thus if the exterior \(L^2\) mass of the relevant Duhamel packet tends to zero, the bounded-frequency \(H^s\) part tends to zero.
The surviving \(H^s\) tail, if it exists, cannot remain in bounded frequency.
It has to escape to \(N_j\to\infty\).

\begin{proposition}[High-frequency alternative]
After the initial tail and compact-core source are removed, a nonzero exterior \(H^s\) source tail produces a high-frequency terminal packet.
\end{proposition}

\begin{proof}
The \(H^s\) norm is the square-summation of dyadic pieces weighted by \(N^s\).
Bounded frequencies vanish by \(L^2\) tightness.
If the full \(H^s\) tail has a positive limsup, some dyadic pieces must remain positive with \(N_j\to\infty\).
Those pieces define the high-frequency terminal packet.
\end{proof}

\section{Fixed Field Scale Versus Dyadic Escape}

The dyadic survivor cannot coexist with a fixed positive Field scale on the same retained Pack+Part branch.
This is the most concrete analytic point in the exterior-source landing.
If Field holds at scale \(\rho>0\) and depth \(M>s\), the packet has finite derivative control through depth \(M\) on finitely many \(\rho\)-scale windows.
For each window \(f\),
\[
  \|P_Nf\|_{L^2}\le C_{\rho,M}N^{-M}\sum_{|\alpha|\le M}\|\partial^\alpha f\|_{L^2}.
\]
Therefore
\[
  N^s\|P_Nf\|_{L^2}\to0.
\]
The finite window sum gives the same limit for the selected packet.

The dyadic survivor asserts a positive lower bound for \(N_j^s\|P_{N_j}f_j\|_{L^2}\) along \(N_j\to\infty\).
The same fixed Field packet cannot satisfy both statements.
Thus, on retained Pack+Part, dyadic escape forces Field failure.

\section{Same-Witness Requirement}

The proof repeatedly says "same witness" because without that requirement the argument would split into unrelated fragments.
Pack failure on one object, Part failure on another object, and Field failure on a third object would not prove anything about the counterexample.
The terminal record \(Q\) keeps the object fixed through the test.

\begin{lemma}[Face failures must belong to the same record]
The conclusion \(\Exit(Q;\Owork)\) is justified only when the failed face is derived for the same terminal record \(Q\) that entered the test.
\end{lemma}

\begin{proof}
Membership is a predicate of \(Q\).
To negate membership, the failed membership condition has to be a condition of \(Q\).
A face failure for a different object would be irrelevant to the record under test.
\end{proof}

This same-witness rule is also what keeps the proof from borrowing a useful estimate from a neighboring branch.
The estimate must return to the record \(Q\).
If it cannot, it is support, not proof.

\section{Translation Of CM Language}

The proof can be read without treating CM as a black box.
The translation is:
\[
\begin{array}{lll}
\Pack_Q & \text{means} & \text{\(Q\) has a positive material body of the same fluid},\\
\Part_{N,Q} & \text{means} & \text{the original equation participates on that body},\\
\Field_{N,r,Q} & \text{means} & \text{a positive-scale continuation readout remains},\\
\Member(Q;\Owork) & \text{means} & \text{all three same-solution structures survive},\\
\Exit(Q;\Owork) & \text{means} & \text{the record fails that membership test}.
\end{array}
\]
With this translation, the argument is not an internal-label proof.
It is a proof that the alleged terminal object either keeps the structures needed for continuation or loses one of the structures needed to be the original solution's terminal member.

\section{What A Referee Must Check}

A referee can check the proof in a finite list.
First, the finite terminal negation supplies a same-solution terminal record.
Second, Pack, Part, and Field are exactly the same-solution structures needed for a member-side terminal branch.
Third, the pass branch invokes the standard \(H^s\) continuation theorem only after those structures survive.
Fourth, the fail branch is the first loss of one of those structures.
Fifth, the whole-space exterior source is not a fourth branch because bounded frequencies vanish and high-frequency survivors fail Field under retained Pack+Part.

The proof fails only if one of those checks fails.
It does not fail because an old positive estimate remains difficult.
The direct positive exterior \(H^s\) theorem would be stronger, but the pass-or-exit proof does not depend on that stronger theorem.

\section{Final Referee Reading}

Read the paper as a terminal-object proof.
The theorem does not say that every bad-looking terminal diagnostic is smooth.
It says that a finite Clay counterexample must produce a same-solution terminal member.
The ordered test decides every such record.
Member records continue.
Nonmember records exit at Pack, Part, or Field.
The whole-space exterior source is decided by the same test after the Duhamel and dyadic reductions.
That is the complete branch structure of the proof.
